\begin{abstract}
    Repository-aware code translation is critical for modernizing legacy systems, enhancing maintainability, and enabling interoperability across diverse programming languages. 
    While recent advances in large language models (LLMs) have improved code translation quality, existing approaches face significant challenges in practical scenarios: insufficient contextual understanding, inflexible prompt designs, and inadequate error correction mechanisms. 
    These limitations severely hinder accurate and efficient translation of complex, real-world code repositories.
            
    To address these challenges, we propose \toolname, a novel multi-agent LLM framework for repository-aware code translation. 
    \toolname systematically decomposes the translation process into specialized subtasks—context retrieval, dynamic prompt construction, and iterative code refinement—each handled by dedicated agents. 
    Our approach leverages retrieval-augmented generation (RAG) for contextual information gathering, employs adaptive prompts tailored to varying repository scenarios, and introduces a reflection-based mechanism for systematic error correction.
            
    We evaluate \toolname on hundreds of Java-C\# translation pairs from six popular open-source projects. 
    Experimental results demonstrate that \toolname significantly outperforms state-of-the-art baselines in both compile and pass rates. 
    Specifically, \toolname achieves up to 55.34\% compile rate and 45.84\% pass rate.
    Comprehensive analysis confirms the robustness and generalizability of \toolname across different LLMs, establishing its effectiveness for real-world repository-aware code translation.
\end{abstract}


\begin{IEEEkeywords}
Multi-Agent, Large Language Model, Code Translation
\end{IEEEkeywords}